% Pracovní list Work/MAT, 1. ročník, 64. hodina.
% Autorské příklady; řešení je v samostatném souboru _reseni.
\documentclass[11pt,a4paper]{article}
\usepackage[T1]{fontenc}
\usepackage[utf8]{inputenc}
\usepackage[czech]{babel}
\usepackage{lmodern,microtype}
\usepackage{amsmath,amssymb,geometry,enumitem,tabularx,booktabs,tikz}
\usetikzlibrary{arrows.meta}
\geometry{top=14mm,bottom=16mm,left=19mm,right=19mm}
\setlength{\parindent}{0pt}\setlength{\parskip}{4pt}
\setlist[enumerate]{leftmargin=1.55em,topsep=2pt,itemsep=2pt}
\pagestyle{plain}
\newcommand{\graf}{%
\begin{tikzpicture}[x=.62cm,y=.62cm,>=Latex]
\draw[step=1,gray!30,very thin] (-5,-5) grid (5,5);
\draw[->] (-5.3,0)--(5.4,0) node[right] {$x$};
\draw[->] (0,-5.3)--(0,5.4) node[above] {$y$};
\foreach \i in {-4,-2,2,4}{\draw (\i,.11)--(\i,-.11) node[below,font=\scriptsize] {$\i$};\draw (.11,\i)--(-.11,\i) node[left,font=\scriptsize] {$\i$};}
\end{tikzpicture}}
\newcommand{\linka}{\par\vspace{2mm}\hrulefill\par}
\begin{document}
{\Large\bfseries Lineární funkce — samostatné procvičování}\hfill{\bfseries 64. hodina}\par
\vspace{1mm}
Jméno a příjmení: \rule{65mm}{.4pt}\hfill Třída: \rule{20mm}{.4pt}\par
\smallskip
\textit{Pracuj s pravítkem. Ke každému grafu napiš předpis přímo do obrázku.
Body zapisuj ve tvaru $[x;y]$. Nepotřebuješ řešit všechny úlohy najednou — postupuj od 1. úlohy.}
\section*{1. Grafy z předpisů}
Do jedné soustavy sestroj grafy $f(x)=x-2$, $g(x)=-2x+1$, $h(x)=\frac12x+2$.
Ke každému grafu urči dvě vhodné dvojice souřadnic. Čáry rozliš různým typem.
\begin{center}\graf\end{center}
\section*{2. Co nám říkají parametry?}
Doplň tabulku podle předpisů z úlohy 1.
\begin{center}
\begin{tabularx}{.94\linewidth}{@{}lcccc@{}}\toprule
Funkce & Směrnice $a$ & Absolutní člen $b$ & Roste / klesá / je konstantní & Bod na ose $y$\\\midrule
$f$ & \rule{11mm}{.35pt} & \rule{11mm}{.35pt} & \rule{32mm}{.35pt} & \rule{17mm}{.35pt}\\[3mm]
$g$ & \rule{11mm}{.35pt} & \rule{11mm}{.35pt} & \rule{32mm}{.35pt} & \rule{17mm}{.35pt}\\[3mm]
$h$ & \rule{11mm}{.35pt} & \rule{11mm}{.35pt} & \rule{32mm}{.35pt} & \rule{17mm}{.35pt}\\\bottomrule
\end{tabularx}
\end{center}
\section*{3. Předpis ze dvou bodů}
Urči předpis $y=ax+b$ procházející a) $A[0;2]$ a $B[2;6]$, b) $C[-2;0]$ a $D[2;-2]$.
Uveď výpočet směrnice i absolutního členu a zkontroluj jeden bod.
\vspace{19mm}
\newpage
\section*{4. Předpis z grafu}
Na obrázku jsou tři různé přímky: \textbf{plná} $p$, \textbf{čárkovaná} $q$ a
\textbf{tečkovaná} $r$. Urči předpisy všech tří funkcí a napiš k nim dva body, které jsi použil.
\begin{center}
\begin{tikzpicture}[x=.85cm,y=.77cm,>=Latex]
\draw[step=1,gray!25,very thin] (-4,-4) grid (4,4);
\draw[->] (-4.35,0)--(4.4,0) node[right] {$x$};
\draw[->] (0,-4.35)--(0,4.4) node[above] {$y$};
\foreach \i in {-4,-3,-2,-1,1,2,3,4}{\draw (\i,.08)--(\i,-.08) node[below,font=\scriptsize] {$\i$};\draw (.08,\i)--(-.08,\i) node[left,font=\scriptsize] {$\i$};}
\begin{scope}\clip (-4,-4) rectangle (4,4);
\draw[very thick] (-4,-9)--(4,7);
\draw[very thick,dashed] (-4,7)--(4,-1);
\draw[very thick,densely dotted] (-4,-1)--(4,3);
\end{scope}
\node[fill=white,inner sep=1pt,font=\small] at (1.7,3.6) {$p$};
\node[fill=white,inner sep=1pt,font=\small] at (3.3,.3) {$q$};
\node[fill=white,inner sep=1pt,font=\small] at (-2.3,-.6) {$r$};
\end{tikzpicture}
\end{center}
$p(x)=$\rule{33mm}{.4pt}\qquad $q(x)=$\rule{33mm}{.4pt}\qquad $r(x)=$\rule{33mm}{.4pt}
\vspace{6mm}\linka
\section*{5. Posouvání grafů}
Do jedné soustavy sestroj $y=x$, $y=x+2$ a $y=x-1$. Popiš, co se změní,
pokud ve funkci $y=x+b$ měníme pouze $b$.
\begin{center}\graf\end{center}
Vysvětlení: \rule{.78\linewidth}{.4pt}
\newpage
\section*{6. Parametr z konečné množiny}
Je dána \emph{rodina funkcí} $y=ax+1$, kde $a\in\{-2;0;1\}$. Napiš všechny
předpisy a sestroj jejich grafy do jedné soustavy.
\begin{center}\graf\end{center}
Kolik různých grafů vznikne? \rule{20mm}{.4pt}\quad Co mají společné? \rule{65mm}{.4pt}
\section*{7. Parametr z intervalu}\textit{Rozšiřující úloha}
Je dána rodina funkcí $y=x+b$, kde $b\in\langle-1;2\rangle$.
Načrtni obě \emph{krajní přímky}, potom několika dalšími přímkami naznač i ostatní.
Kolik různých grafů této rodině odpovídá? Je to jediná funkce?
\begin{center}\graf\end{center}
Odpověď: \rule{.82\linewidth}{.4pt}
\newpage
\section*{8. Omezený definiční obor}\textit{Rozšiřující úloha}
Nakresli graf funkce $f(x)=-x+2$ jen pro $x\in\langle-2;1)$.
Vyznač oba koncové body tak, aby bylo poznat, který do grafu patří.
Urči $D(f)$ a $H(f)$.
\begin{center}\graf\end{center}
$D(f)=$\rule{50mm}{.4pt}\qquad $H(f)=$\rule{50mm}{.4pt}
\section*{9. Dva body a posun}\textit{Dobrovolná rezerva}
Urči předpis přímky procházející body $A[1;3]$ a $B[3;7]$. Porovnej jej s přímkou
$y=2x-1$. Jsou rovnoběžné? Pokud ano, o kolik je jedna výše?
\par\vspace{27mm}
\noindent\textit{Po práci si zkontroluj:} Rozlišoval(a) jsem argument $x$ od parametru?
Přidělil(a) jsem každému grafu předpis? Zohlednil(a) jsem otevřený krajní bod?
\end{document}
